\section{BIG-Bench Hard：从平均分到任务级异质性}

单个 arithmetic demo 可能只是 cherry-pick，因此课件转向 BIG-Bench Hard（BBH）：从 BIG-Bench 中选出当时模型低于 average human rater 的困难任务，再系统比较 answer-only 与 CoT。这个设计把“是否能做复杂推理”拆成导航、排序、逻辑和状态追踪等多种 task mechanics。

\subsection{Benchmark 怎样定义“hard”}

BBH 的难度来自构造条件，而不是永久属性：被选任务在原 BIG-Bench 结果中没有模型超过 average human。Prompt 包含 task description、worked examples、options 和 reasoning trace；因此评估的是模型能否从自然语言描述与 exemplars 中推断 algorithm。

\lecturefigure{slide-24.jpg}{BIG-Bench Hard 的任务选择与 CoT prompt 结构}{官方课件第 24 页；Suzgun et al., 2022}

\begin{knowledgebox}{读图：两个例子测量不同的中间状态}
Navigation 要连续更新方向与坐标，最后判断是否回到原点；word sorting 要比较首字母、处理相同前缀并维护输出顺序。两项任务对人类都不神秘，却要求模型在多个 token 之间保持一致状态。CoT 的价值是把状态更新显式写出来，而不是只在最后猜 yes/no 或单个列表。
\end{knowledgebox}

\subsection{平均分之外，还要看多少任务跨过人类线}

若只报告 aggregate accuracy，少数高收益任务可能掩盖大量失败。课件同时报告平均 performance 和“超过 average human rater 的任务比例”，并用红蓝任务图展示异质性。两种 metric 分别回答“总体提高多少”和“能力覆盖面扩大多少”。

\lecturefigure{slide-25.jpg}{BBH 结果总览：平均性能、超过人类的任务数与任务级地图}{官方课件第 25 页；Suzgun et al., 2022}

\begin{knowledgebox}{读图：先理解 subset bias，再理解增益}
BBH 本来就选择旧模型低于人类的任务，所以 answer-only baseline 较低是 by construction。加入 CoT 后，davinci-002 的平均分和超过 average-human 的任务数都上升；下方红蓝图说明增益并不均匀。结论应是 CoT 解锁了多数但非全部困难任务，而不是“模型全面超过人类”。
\end{knowledgebox}

\teachervoice{课堂逐项解释两个 metric，并指出 average human 约为 67、best human 更高。讲者没有把“多数任务超过 average human”写成通用人类水平，而是把它限定在 BBH 的题集、prompt 与评分方式内。}

\subsection{Scaling：CoT 的正 delta 有自己的阈值}

下一张图把最大模型结果拆回多个 scale。若 CoT 只是在 prompt 中增加信息，所有模型都应得到类似相对收益；实际曲线却显示，小模型不一定获得正 delta，直到 davinci-002 / PaLM 62B 附近才在 aggregate 上稳定受益。这个 scale sweep 是方法论上的关键，因为它把“最大模型上的成功案例”升级为对 intervention 与 base capability 交互作用的检验。

\lecturefigure{slide-26.jpg}{BBH 上的 CoT scaling：正收益需要足够模型规模}{官方课件第 26 页；Suzgun et al., 2022}

\begin{knowledgebox}{读图：delta 比绝对分更能识别方法门槛}
每个 panel 中比较同一模型的 no-CoT 与 CoT，而不是只比较不同模型绝对分数。小模型的 reasoning trace 可能引入额外错误；大模型能利用 trace 分解任务，才出现稳定正差。图中的阈值是特定 model family 与 task aggregate 的经验位置，不应被写成所有 CoT 系统的固定 62B 规则。
\end{knowledgebox}

\subsection{Emergence：No-CoT 仍平，CoT 曲线开始上升}

最强的课堂叙事来自对照：answer-only performance 在被测规模中仍近乎 flat，CoT 却让较大模型显著上升。这里“解锁”指的是 interface 让已有但不可直接调用的组合能力进入 metric，而不是证明 pretraining 后模型内部完全没有相关知识。

\lecturefigure{slide-27.jpg}{BBH 上的 CoT emergence：answer-only 平坦，CoT 打开性能增长}{官方课件第 27 页；Suzgun et al., 2022}

\begin{knowledgebox}{读图：Prompt 可以改变我们看到的能力边界}
若只测试 no-CoT，研究者可能得出“继续扩大也没有进展”；加入 reasoning exemplars 后，同一模型族出现规模相关增益。这说明 capability evaluation 不能与 elicitation method 分离。图支持“CoT 在此设置中提升可测能力”，但不证明所有 flat curve 都能被更好 prompt 解锁。
\end{knowledgebox}

\begin{importantbox}{评估模型能力时要同时报告 elicitation budget}
模型、prompt、few-shot examples、generated-token budget、sampling strategy 和 external tools 共同定义一次能力测试。只报告 parameter count 与 final accuracy，会把 interface failure 误写成 capability absence，也会把昂贵 inference-time search 隐藏成“模型本身更聪明”。
\end{importantbox}

\subsection{本章小结}

BBH 将 CoT 从单例扩展到困难任务集合。平均分、跨人类任务比例和任务级地图共同显示异质收益；scale sweep 又说明 CoT 本身是 emergent prompting technique。能力曲线因此必须与 elicitation method 一起解释。

